\chapter{La langue des épitaphes}\label{ch:langue}

Les épitaphes datées offrent un échantillon daté du latin écrit à Rome dans la seconde moitié du \textsc{iv}\textsuperscript{e} siècle. Elles ne donnent pas pour autant un accès direct à la langue parlée. Le texte gravé passe par un rédacteur, un modèle de formulaire et un lapicide ; le texte édité passe ensuite par les normalisations de l’éditeur et de la base. Ce chapitre décrit ce que le texte transmis permet d’observer, en séparant les faits de graphie, les abréviations et les choix morphologiques qui relèvent du développement éditorial.

\section{Graphies : ce que signalent les corrections éditoriales}

Le texte de l’EDCS est normalisé : la lettre attendue figure dans le texte, la lettre déclarée gravée est conservée dans la notation \emph{<attendu=gravé>}, rendue ici par les demi-crochets ⸢ ⸣ et l’apparat de chaque notice. L’état 2022 de la base signale ainsi \st{nCorrOcc} corrections dans \st{nCorrNotices} notices de l’ensemble A. Ces données ne constituent pas une transcription diplomatique : elles reflètent ce que les éditeurs ont jugé utile de signaler, et une graphie non standard qu’un éditeur aurait conservée sans la corriger n’y figure pas.

Leur typologie est néanmoins nette (tableau~\ref{tab:graphies}). Quatre phénomènes dominent :

\begin{itemize}
\item \emph{le bêtacisme}, échange de B et de V (81 occurrences) : \emph{B} pour \emph{v} dans 72 cas (\emph{⸢b⸣ixit}, \emph{iu⸢b⸣enis}), \emph{V} pour \emph{b} dans 9 ;
\item \emph{la confusion de E et de I} (60 occurrences), dans les deux sens (\emph{mens⸢e⸣s} gravé MENSIS, \emph{⸢i⸣n} gravé EN) ;
\item \emph{la confusion de O et de U} (46 occurrences), notamment dans la désinence d’\emph{annos}, gravée ANNVS (28 occurrences dans 26 notices), et dans les formes \emph{Theudosius} et \emph{Neuterius} des noms consulaires ;
\item \emph{les graphies du Y} (18 occurrences), rendu par I ou par V dans les noms d’origine grecque (\emph{Eut⸢y⸣chius}).
\end{itemize}

Viennent ensuite C pour G (13), N pour M devant labiale ou en finale (12), S pour X (6) et E pour A (5). Ces phénomènes sont ceux que décrivent les grammaires du latin tardif ; ils ne sont pas propres aux inscriptions chrétiennes.\footnote{V.~Väänänen, \emph{Introduction au latin vulgaire}, 3\textsuperscript{e} éd., Paris, 1981 ; J.~N.~Adams, \emph{Social Variation and the Latin Language}, Cambridge, 2013. Notices bibliographiques vérifiées ; aucun renvoi de page n’est donné, faute de consultation des passages.} Une partie des corrections relève par ailleurs de la lecture plutôt que de la langue : les dix-huit cas où un I est lu à la place d’un T attendu peuvent tenir à une barre horizontale effacée ou mal gravée, ce qu’une photographie seule permettrait de décider.

% généré par scripts/04_generated.py
\begin{longtable}{@{}l l r l p{0.45\linewidth}@{}}
\caption{Corrections éditoriales signalées par l’EDCS (état 2022, ensemble A) : lettre attendue, lettre déclarée gravée, occurrences, phénomène, notices. Correspondances attestées au moins deux fois.}\label{tab:graphies}\\
\toprule Attendu & Gravé & Occ. & Phénomène & Notices\\\midrule\endfirsthead\toprule Attendu & Gravé & Occ. & Phénomène & Notices\\\midrule\endhead\bottomrule\endfoot
\lat{v} & B & 72 & bêtacisme & {\footnotesize 4, 6, 22, 24, 26, 34, 40, 47, 56, 63, 72, 83, 89, 90, 91, 106, 112, 113, 118, 119, 120, 124, 144, 172, 184, 194, 207, 208, 215, 217, 234, 235, 261, 272, 281, 301, 313, 332, 347, 352, 365, 366, 374, 402, 407, 408, 417, 440, 441, 444, 468, 475} \\
\lat{o} & V & 44 & confusion o/u & {\footnotesize 32, 54, 55, 69, 94, 104, 130, 161, 191, 193, 209, 225, 275, 281, 288, 289, 299, 320, 322, 323, 327, 336, 346, 347, 353, 359, 365, 374, 375, 376, 379, 400, 419, 421, 427, 431, 443, 466, 470, 473, 475, 478} \\
\lat{e} & I & 43 & confusion e/i & {\footnotesize 4, 25, 28, 34, 55, 60, 80, 84, 94, 124, 151, 157, 191, 193, 260, 261, 287, 306, 337, 347, 353, 359, 365, 366, 367, 381, 389, 390, 395, 411, 423, 431, 441, 443, 446} \\
\lat{t} & I & 18 & i pour t (lecture) & {\footnotesize 37, 63, 132, 136, 153, 164, 186, 231, 245, 296, 306, 308, 333, 371, 433} \\
\lat{i} & E & 17 & confusion e/i & {\footnotesize 27, 40, 50, 54, 70, 76, 100, 111, 125, 306, 353, 364, 368, 388, 391, 443} \\
\lat{y} & I & 13 & graphie de y & {\footnotesize 1, 145, 172, 233, 238, 276, 308, 352, 391, 417, 418, 419, 424} \\
\lat{g} & C & 13 & c pour g & {\footnotesize 19, 35, 99, 201, 221, 231, 304, 306, 352, 353, 354, 359, 366} \\
\lat{m} & N & 12 & n pour m & {\footnotesize 23, 36, 46, 119, 126, 184, 217, 275, 365, 419, 424} \\
\lat{b} & V & 9 & bêtacisme & {\footnotesize 85, 123, 125, 276, 340, 357, 380, 382, 470} \\
\lat{V} & B & 7 &  & {\footnotesize 55, 90, 97, 128, 416} \\
\lat{x} & S & 6 & s pour x & {\footnotesize 13, 130, 132, 272, 302, 304} \\
\lat{l} & I & 6 &  & {\footnotesize 186, 266, 287, 333, 359} \\
\lat{a} & E & 5 & e pour a & {\footnotesize 28, 214, 279, 466} \\
\lat{y} & V & 5 & graphie de y & {\footnotesize 223, 239, 338, 339, 340} \\
\lat{q} & O & 4 &  & {\footnotesize 111, 231, 388, 418} \\
\lat{o} & VM & 4 &  & {\footnotesize 291, 402} \\
\lat{c} & Q & 3 & q pour c & {\footnotesize 1, 184, 287} \\
\lat{K} & C & 3 &  & {\footnotesize 9, 128, 279} \\
\lat{c} & QV & 3 &  & {\footnotesize 59, 404, 405} \\
\lat{t} & D & 3 &  & {\footnotesize 182, 319, 402} \\
\lat{a} & O & 3 &  & {\footnotesize 448, 449, 451} \\
\lat{i} & V & 2 &  & {\footnotesize 44, 157} \\
\lat{vi} & BE & 2 &  & {\footnotesize 66, 458} \\
\lat{n} & V & 2 &  & {\footnotesize 85, 309} \\
\lat{n} & M & 2 &  & {\footnotesize 121, 150} \\
\lat{XV} & VX & 2 &  & {\footnotesize 128, 198} \\
\lat{u} & O & 2 & confusion o/u & {\footnotesize 143, 299} \\
\lat{i} & O & 2 &  & {\footnotesize 156, 275} \\
\lat{e} & Y & 2 &  & {\footnotesize 198, 303} \\
\lat{o} & I & 2 &  & {\footnotesize 231, 402} \\
\lat{o} & U & 2 &  & {\footnotesize 260, 282} \\
\lat{r} & B & 2 &  & {\footnotesize 308, 371} \\
\lat{r} & I & 2 &  & {\footnotesize 446, 464} \\
\end{longtable}
\noindent{\small Correspondances attestées une seule fois (54) : d/T (7); Fl/IT (7); ae/I (26); v/R (34); c/I (37); ag/VA (40); st/TS (60); C/K (62); tt/DD (72); i/U (115); te/S (131); o/EM (133); d/P (135); V/Y (142); s/RO (156); q/C (184); ibus/ATV (193); ba/PR (201); us/OM (204); e/A (205); o/A (205); G/C (211); h/A (215); F/E (219); ch/T (223); v/I (231); y/U (232); r/D (269); d/O (287); i/B (287); it/TI (296); b/P (299); hr/RH (306); te/I (306); is/SI (309); L/I (318); e/C (319); c/E (319); um/ON (320); t/S (336); i/C (340); u/Y (349); x/CS (353); i/Y (368); ly/IV (371); ri/K (371); ae/IS (381); r/S (404); s/R (406); rum/S (416); S/E (441); i/A (470); e/F (473); Fl/TI (473).}


Le corpus ne permet pas, en général, de mesurer une évolution chronologique de ces graphies : les corrections sont dispersées, et leur signalement dépend de l’éditeur de chaque volume des ICUR autant que du texte gravé. Une correspondance fait exception par sa répartition. La graphie ANNVS pour \emph{annos} apparaît dans 2 notices sur 55 en 340–359, 4 sur 123 en 360–374, 4 sur 139 en 375–389, puis dans 16 notices sur 160 en 390–403 (ensemble A, hors \nnum{98}). La concentration dans la dernière période est nette ; elle peut traduire l’extension d’une graphie phonétique de la finale, mais elle pourrait aussi tenir à l’atelier ou au secteur d’origine des textes, que l’on ne connaît que pour une partie d’entre eux. Elle mérite un contrôle sur les photographies. L’index des graphies (chapitre~\ref{ch:index}) et le répertoire du chapitre~\ref{ch:rep_graphies} donnent pour chaque correspondance la liste des notices.

\section{Abréviations}

La datation consulaire fournit l’essentiel des abréviations (tableau~\ref{tab:abrev}) : \emph{conss(ulibus)}, \emph{cons(ulibus)}, \emph{vv(iris) cc(larissimis)}, \emph{dd(ominis) nn(ostris)}, \emph{Aug(usto)}, \emph{Fl(avio)}. Le redoublement de la dernière consonne pour marquer le pluriel (\emph{conss.}, \emph{Augg.}, \emph{vv. cc.}) est fréquent. Viennent ensuite les éléments de l’âge (\emph{ann(os)}, \emph{m(enses)}, \emph{d(ies)}) et de la date (\emph{Kal(endas)}, \emph{dep(ositus)}, \emph{d(e)p(ositus)}). Dans l’ensemble C, la formule consulaire n’est écrite \emph{consulibus} en toutes lettres que dans sept textes (\nnums{7}, 22, 155, 199, 215, 260 et 265) ; la forme abrégée est la norme.

% généré par scripts/04_generated.py
\begin{table}[tbp]\centering\small\caption{Les abréviations les plus fréquentes de l’ensemble A : forme abrégée (lettres gravées), exemple de développement éditorial, occurrences.}\label{tab:abrev}
\begin{tabular}{@{}l l r@{}}\toprule Abrégé & Exemple & Occ.\\\midrule d & \lat{d(ies)} & 133 \\
kal & \lat{Kal(endas)} & 124 \\
conss & \lat{conss(ulibus)} & 119 \\
m & \lat{m(enses)} & 100 \\
aug & \lat{Aug(usto)} & 98 \\
ann & \lat{ann(os)} & 88 \\
fl & \lat{Fl(avio)} & 86 \\
cons & \lat{cons(ulibus)} & 76 \\
n & \lat{n(ostro)} & 68 \\
v & \lat{v(iro)} & 61 \\
c & \lat{c(larissimo)} & 60 \\
p & \lat{p(lus)} & 60 \\
que & \lat{qu(a)e} & 56 \\
dep & \lat{dep(ositus)} & 47 \\
cc & \lat{cc(larissimis)} & 37 \\
vv & \lat{vv(iris)} & 35 \\
dd & \lat{dd(ominis)} & 28 \\
nn & \lat{nn(ostris)} & 23 \\
dp & \lat{d(e)p(osita)} & 23 \\
an & \lat{an(nos)} & 19 \\
q & \lat{q(uae)} & 19 \\
non & \lat{Non(as)} & 17 \\
oct & \lat{Oct(obres)} & 16 \\\bottomrule\end{tabular}\hspace{1.5em}\begin{tabular}{@{}l l r@{}}\toprule Abrégé & Exemple & Occ.\\\midrule coss & \lat{co(n)ss(ulibus)} & 15 \\
vix & \lat{vix(it)} & 15 \\
dec & \lat{Dec(embres)} & 14 \\
sept & \lat{Sept(embres)} & 12 \\
ian & \lat{Ian(uarias)} & 12 \\
id & \lat{Id(us)} & 12 \\
iul & \lat{Iul(ias)} & 12 \\
feb & \lat{Feb(ruarias)} & 12 \\
s & \lat{s(upra)} & 11 \\
augg & \lat{Augg(ustis)} & 10 \\
k & \lat{K(alendas)} & 10 \\
mai & \lat{Mai(as)} & 10 \\
f & \lat{f(emina)} & 9 \\
cos & \lat{co(n)s(ulibus)} & 9 \\
con & \lat{con(sulibus)} & 8 \\
teodosio & \lat{T(h)eodosio} & 8 \\
aur & \lat{Aur(elio)} & 7 \\
mart & \lat{Mart(ias)} & 7 \\
mens & \lat{mens(es)} & 7 \\
cl & \lat{Cl(audio)} & 7 \\
nov & \lat{Nov(embres)} & 7 \\
pos & \lat{pos(t)} & 6 \\\bottomrule\end{tabular}\end{table}


\section{Morphologie : ce qui appartient à l’éditeur}

Le développement des abréviations est un acte éditorial. Il engage parfois un choix morphologique que la pierre ne tranche pas. L’exemple le plus net est l’expression de la durée de vie. Le texte normalisé de l’ensemble C contient la forme \emph{annos} dans 220 textes ; mais elle n’est écrite en toutes lettres que dans 45 d’entre eux. Il faut y ajouter les textes où la pierre porte ANNVS, que l’éditeur corrige en \emph{ann⸢o⸣s} : la désinence est alors gravée, sous une forme non standard. Dans les autres, la pierre porte \emph{ann.}, \emph{an.} ou un sigle, et l’accusatif est une décision de l’éditeur. À l’inverse, l’ablatif \emph{annis}, attesté dans 40 textes, est écrit en toutes lettres dans 29. Il en va de même pour \emph{menses} (103 textes, dont 15 en toutes lettres) et \emph{dies} (90 textes, dont 31 en toutes lettres).

Sur les seules formes écrites, l’accusatif de durée (\emph{vixit annos}) et l’ablatif (\emph{vixit annis}) coexistent dans toutes les périodes : 8 contre 8 textes en 340–359, 9 contre 7 en 390–403. Les effectifs sont trop faibles pour décrire une tendance, mais ils suffisent à montrer qu’une étude de la syntaxe de la durée fondée sur le texte développé mesurerait d’abord les habitudes des éditeurs.

La même prudence s’applique aux formes de la déposition. Le texte normalisé donne \emph{depositus} dans 122 textes et \emph{deposita} dans 109, mais ces formes ne sont écrites en toutes lettres que dans 36 et 47 textes ; ailleurs, \emph{dep.} ou \emph{d. p.} est développé selon le sexe du défunt, lorsque celui-ci est connu par ailleurs, ou selon l’usage de l’éditeur. Le substantif \emph{depositio} n’est écrit en toutes lettres qu’une fois (\nnum{448}).

\section{Registre et composition}

La très grande majorité des textes relève d’un formulaire bref, où la langue se réduit à quelques syntagmes figés. Une dizaine de textes s’en écartent par la versification ou par une prose plus élaborée : les épitaphes métriques de Felicitas (\nnum{470}), d’Eutychius (\nnum{352}) ou de Maximilla (\nnum{306}), les textes du \nnum{232} et du \nnum{73}. Ils sont étudiés pour leur contenu aux chapitres~\ref{ch:dieu_christ} à \ref{ch:textes_decisifs}. Leur langue est celle de la poésie funéraire latine, avec ses emprunts à la tradition classique (\emph{funus crudele} au \nnum{352} ; \emph{lumina}, \emph{manes} au \nnum{470}) ; la présence d’un vocabulaire chrétien n’y exclut pas le recours aux formules païennes de la consolation.
